\section{Background}\label{sec:prelim}
\subsection{Masked Diffusion Models}
A continuous-time discrete diffusion model consists of a forward noising process and a backward denoising process, both continuous-time Markov chains. In the forward process, training samples are progressively corrupted into pure noise. The model learns the corresponding backward process that inverts this corruption, then generates new samples by running the backward process from noise.

\textbf{Continuous-time Markov chain.} We consider a discrete state space $\mathcal{X}$. A continuous-time Markov chain (CTMC) is a process $\vx_t$ on $\mathcal X$ with ${t \in [0, 1]}$, starting from an initial data distribution $q_0(\vx_0)$. A CTMC is characterized by a time-dependent transition rate matrix $Q_t \in \mathbb{R}^{|\mathcal{X}| \times |\mathcal{X}|}$. For distinct states $\vx_t, \vy \in \mathcal{X}$, $Q_t(\vx_t, \vy) \ge 0$ defines the instantaneous transition rate from $\vx_t$ to $\vy$, i.e., over an infinitesimal time interval $[t,t+\Delta t]$ the transition probability is given by
\begin{equation}\label{eq:infinitesimal_transition}
    q_{t+\Delta t|t}(\vy|\vx_t) = \delta(\vx_t, \vy) + Q_t(\vx_t, \vy) \Delta t + o(\Delta t),
\end{equation}
where $\delta$ is the Kronecker delta.
Equivalently, the marginal distribution $q_t\in\Delta(\mathcal X)$ evolves according to the Kolmogorov ODE $dq_t/dt = q_t Q_t$.
Note that the diagonal entries of $Q_t$ should satisfy $Q_t(\vx_t, \vx_t) = - \sum_{\vx_t \neq \vy} Q_t(\vx_t, \vy)$ to ensure that $q_t$ sums to $1$. 

\textbf{The masking and unmasking processes.} In the forward process of MDMs with law $q_t$, the transition rate matrix is commonly parameterized as $Q_t = \sigma(t)Q$, where $\sigma(t)$ is a scalar noise schedule and $Q$ is a constant rate matrix determined by the model design \citep{campbell2022continuous}. For MDMs, a specific $Q$ is adopted so that over an infinitesimal interval $[t,t+\Delta t]$, the token-level transition probabilities are
\begin{equation*}
    q_{t+\Delta t|t}(y|x)=
\begin{cases}
\sigma(t)\,\Delta t + o(\Delta t), & y=\mask,\\
1-\sigma(t)\,\Delta t - o(\Delta t), & y=x\in\mathcal V,\\
0, & \text{otherwise},
\end{cases}
\end{equation*}
with a state space $\mathcal V\cup\{\mask\}$. Letting $\bar\sigma(t)=\int_0^t\sigma(s)\,ds$, the distribution of a token at time $t$ conditioned on the initial state $x$ is given by
\begin{equation*}
\label{eq:forward_seq}
    q_{t|0}(y|x)
=
e^{-\bar\sigma(t)}\mathbf 1[y=x]
+
\bigl(1-e^{-\bar\sigma(t)}\bigr)\mathbf 1[y=\mask].
\end{equation*}
Due to token-wise independence, the forward process for a length-$L$ sequence is the product CTMC on $(\mathcal V\cup\{\mask\})^L$, with the joint distribution
\begin{equation*}
\label{eq:forward_distribution}
        q_{t|0}(\bm x_t|\bm x_0)=\prod_{i=1}^L q_{t|0}(x_t^i|x_0^i).
\end{equation*}

The model aims to learn the time-reversal of the forward process $q_t$, denoted by $p_t^\theta$: starting from an all-mask terminal state, we randomly unmask each position to a token in the vocabulary. We provide mathematical details in Appendix~\ref{sec:literature-appendix}.


\subsection{Related Works on MDM Inference}\label{subsec:confidence}

\textbf{Confidence-based unmasking.}
A common alternative to uniform unmasking in masked diffusion models (MDMs) is to prioritize positions where the model exhibits higher predictive confidence. Concretely, given the conditional token distribution
$p_i(v) = p_t^\theta(x_{t-\Delta t}^i = v \mid \bm x_t)$
at position $i$, a confidence score $c_i$ is computed and used to guide the unmasking order. Popular choices include:
\begin{itemize}[noitemsep, topsep=0pt]
    \item {Negative entropy} \citep{ye2025dream7bdiffusionlarge}, $c_i = \sum_v p_i(v)\log p_i(v)$;
    \item {Top-1 probability} \citep{nie2025scalingmaskeddiffusionmodels}, $c_i = \max_v p_i(v)$;
    \item {Top-2 margin} \citep{kim2025train}, $c_i = p_i(v_1) - p_i(v_2)$, where $v_1, v_2$ are the two most likely tokens at $i$.
\end{itemize}
Given these scores, unmasking strategies are typically either \emph{greedy}, selecting the top-$N$ positions with highest confidence, or \emph{randomized}, for example by perturbing confidence scores before top-$N$ selection \citep{kim2025train} or sampling positions proportionally via a softmax transformation \citep{ye2025dream7bdiffusionlarge}. While simple and effective in practice, these approaches are largely heuristic and do not arise from an explicit optimization objective. Beyond fixed-size selection, \citet{benhamu2025eb} propose an adaptive sampler that dynamically chooses both which and how many tokens to unmask, justified by a KL-based entropy bound on joint dependence error.

\textbf{Trainable inference methods.}
Beyond fixed heuristics, several works treat the unmasking order as a trainable policy. We defer these discussions to Appendix~\ref{sec:literature-appendix}. 

\iffalse
\textbf{Theoretical perspectives and limitations.}
Recent theoretical work has begun to clarify why inference strategies play a central role in discrete and masked diffusion models. Formalizations of discrete diffusion as continuous- or discrete-time Markov processes show that commonly used approximations, such as $\tau$-leaping, can introduce systematic bias that compounds in high-dimensional discrete spaces \citep{austin2021structured, campbell2022continuous}. From a Gibbs sampling perspective, diffusion inference can be viewed as a sequence of conditional updates, where the unmasking order acts as a scan rule that critically affects convergence and stability \citep{hayakawa2025demystifying, peng2025path}. 

These analyses also help explain the limitations of confidence-based heuristics. Greedy unmasking based on local confidence resembles coordinate descent in an implicit energy landscape, making it highly sensitive to early-stage errors that can propagate and trap the sampler in suboptimal regions \citep{kholkin2025sampling, chen2025beyond}. While learning-based approaches partially address these issues, they often optimize over local transitions or trajectories rather than globally optimal paths through the diffusion process.

\textbf{Our perspective.}
Building on these insights, we adopt a more foundational viewpoint: rather than treating confidence heuristics or learned policies as ad hoc design choices, we ask what inference objective they implicitly approximate. This perspective allows us to (i) derive confidence-based unmasking from a directly trainable global objective, and (ii) explain when and why such heuristics become suboptimal under certain data or model regimes.
\fi

\iffalse

\textbf{Confidence-based methods.} As an alternative to uniform unmasking, existing works propose to prioritize unmasking tokens where the MDM's prediction confidence is higher. In these methods, we first compute MDM's confidence $c_i$ at each token position $i$ using the predicted token probabilities $p_i(v)=p_t^\theta(x_{t-\Delta t}^{i}=v|\bm x_t)$. Common metrics of confidence include
\begin{itemize}[noitemsep, topsep=0pt]
    \item negative entropy \citep{ye2025dream7bdiffusionlarge}, where $c_i=\sum_v p_i(v)\log p_i(v)$;
    \item top-1 probability \citep{nie2025scalingmaskeddiffusionmodels}, where $c_i=\max_v p_i(v)$;
    \item top-2 probability margin \citep{kim2025train}, where $c_i=p_i(v_1)-p_i(v_2)$ and $v_1, v_2$ denotes the 2 tokens with highest probabilities among $\mathcal V$.
\end{itemize}
With the computed confidence, one can design a family of unmasking strategies that are either \textit{greedy}---directly select top-$N$ indices with the highest confidence to unmask, or \textit{randomized}---adding noise to the confidence score before top-$N$ \citep{kim2025train} or converting into unmasking probabilities by softmax function \citep{ye2025dream7bdiffusionlarge}.

\textbf{Trainable methods.} 

Several works consider the ordering policy as a separate trainable entity \citep{peng2025path,jazbec2025learning} or construct update rules that allow the backwards process to learn policies that are more globally optimal, such as updating coherent groups of variables that share semantic dependencies \citep{meshchaninov2025guided} or using informed corrective steps to modify unmasked tokens \citep{zhao2024informed}. However, one weakness of these works is the lack of generalizable training framework that would illustrate the specific training objective that an optimal policy should optimize.


\textbf{Theory and understanding.} 

However, confidence based methods have several problems. First, most empirical methods are purely heuristics without an explicit training objective \citep{kim2025train} or learn a policy through reinforcement learning \citep{jazbec2025learning}. Second, the obtained policies optimize over local paths rather than global paths, which can lead to early consistency errors that propagate downstream \citep{chen2025beyond, hayakawa2025demystifying}. Recent theoretical work has established that inference in discrete and masked diffusion models is fundamentally governed by the sampling policy rather than the denoiser alone. Formalizations of discrete diffusion as discrete-time or continuous-time Markov processes show that commonly used approximations such as $\tau$-leaping can accumulate bias in high-dimensional discrete spaces \citep{austin2021structured, campbell2022continuous}. Interpreting diffusion sampling as Gibbs-style updates further reveals that the update order functions as a scan rule, which critically affects convergence and stability, motivating structured or planned update schedules \cite{hayakawa2025demystifying, peng2025path}. Energy-based perspectives partially explain why greedy confidence-based unmasking behaves like coordinate descent and is sensitive to early-stage errors, which can propagate and trap the sampler in suboptimal regions \cite{kholkin2025sampling, chen2025beyond}. We build on these ideas through a more foundational approach to understand how confidence methods can be derived from a directly trainable objective and also explain the suboptimality seen in certain data regimes.
\fi

% \sijin{Old version: Several recent works focus on selecting good sampling trajectories and discourage bad ones. This is a vital part of Seed Diffusion \citep{song2025seeddiffusionlargescalediffusion}. There are many ways to learn the conditional unmasking probabilities of the backwards process. However, constructing a sampling policy given these conditional probabilities can be difficult. In the following sections, we produce a framework that both formalizes the optimization problem of producing an optimal policy and show how previous heuristics arise from simplifications of the objective function.}